\chapter{Dry Orgasm}

\section*{Definition and Overview}
Dry orgasm is the experience of orgasm without the release of semen. In men, a normal ejaculation involves the forceful release of semen through the penis, but with a dry orgasm the pleasurable climax occurs without ejaculate. The condition is sometimes called retrograde ejaculation, in which semen is propelled backwards into the bladder rather than out through the urethra, or it may involve orgasm with reduced or absent seminal fluid for other reasons. This chapter explains the causes of dry orgasm, how it is diagnosed, and the treatment options, as well as the important distinction between dry orgasm and ejaculation problems that affect fertility.

\section*{Epidemiology}
Dry orgasm is a relatively common problem, particularly in men who have had surgery on the prostate or bladder neck, and its prevalence increases with age. Retrograde ejaculation affects a substantial proportion of men who have undergone transurethral resection of the prostate or bladder neck surgery. Many men are unaware that their dry orgasm is retrograde because the semen is simply redirected into the bladder and later passed in urine. The condition is often first identified when a couple seeks help for infertility.

\section*{Aetiology and Pathophysiology}
Ejaculation requires coordinated contraction of the muscles of the pelvic floor and closure of the bladder neck, which directs semen through the urethra. Dry orgasm occurs when this coordination fails or when seminal fluid production is reduced.
\begin{itemize}
    \item \textbf{Retrograde ejaculation:} the bladder neck fails to close during ejaculation, so semen flows backwards into the bladder; causes include prostate surgery, bladder neck surgery, diabetes-related nerve damage, and some medications
    \item \textbf{Prostate surgery:} transurethral resection of the prostate and other prostate operations commonly cause retrograde ejaculation
    \item \textbf{Diabetes:} long-standing diabetes damages the autonomic nerves that control the bladder neck
    \item \textbf{Medications:} alpha-blockers used for prostate problems, some blood pressure medicines, and antidepressants can reduce or reverse ejaculation
    \item \textbf{Neurological conditions:} spinal cord injury, multiple sclerosis, and other nerve disorders
    \item \textbf{Reduced seminal fluid:} low testosterone, obstruction of the ejaculatory ducts, or removal of the seminal vesicles or prostate
    \item \textbf{Ageing:} the volume of ejaculate naturally declines with age
\end{itemize}

\section*{Clinical Features}
The key feature is an orgasm without semen, which the man may or may not notice depending on his expectations.
\begin{itemize}
    \item The sensation of orgasm is normal, but no ejaculate is released
    \item Urine passed after orgasm may appear cloudy or milky, particularly with retrograde ejaculation
    \item Reduced volume of ejaculate may also be described
    \item Fertility problems when a couple is trying to conceive
    \item Possible discomfort or pain with ejaculation, depending on the cause
\end{itemize}

\section*{Differential Diagnosis}
Dry orgasm should be distinguished from other ejaculatory disorders.
\begin{itemize}
    \item Premature ejaculation, in which ejaculation occurs too quickly
    \item Delayed ejaculation or anejaculation, in which orgasm and ejaculation are difficult or absent
    \item Low ejaculate volume from hormonal or ductal causes
    \item Erectile dysfunction, which is a separate condition affecting arousal
    \item Painful ejaculation, which may indicate infection or obstruction
\end{itemize}

\section*{When to Seek Medical Care}
Men who notice a change in ejaculation, especially after surgery or with new medications, should mention it to their doctor. Medical review is particularly important if the man wishes to father children, if there is blood in the semen or urine, or if the change is accompanied by pain or other urinary symptoms.

\textbf{Call 000 (or local emergency services) for:}
\begin{itemize}
    \item Sudden severe pain in the testicles, groin, or abdomen
    \item Blood in the urine accompanied by pain or inability to pass urine
    \item Sudden inability to pass urine at all
\end{itemize}

\section*{Diagnosis and Investigations}
Diagnosis is usually straightforward and begins with the history.
\begin{itemize}
    \item \textbf{History:} details of ejaculation, surgery, medications, diabetes, and any fertility concerns
    \item \textbf{Urine examination:} checking for semen in urine passed after orgasm confirms retrograde ejaculation
    \item \textbf{Blood tests:} testosterone levels and blood sugar checks where appropriate
    \item \textbf{Semen analysis:} for men concerned about fertility, to assess sperm numbers in any available ejaculate
    \item \textbf{Referral:} to a urologist for further assessment in complex cases
\end{itemize}

\section*{Management}
Treatment depends on the cause and whether fertility is a goal.
\begin{itemize}
    \item \textbf{Reassurance:} dry orgasm is harmless in itself, and many men need no treatment once they understand why it happens
    \item \textbf{Medication adjustment:} where a medicine is responsible, a doctor may change the dose or the drug
    \item \textbf{Medication for retrograde ejaculation:} some medicines can help close the bladder neck and restore normal ejaculation, though success varies
    \item \textbf{Fertility treatment:} where conception is desired, sperm can often be recovered from urine after orgasm and used in assisted reproduction
    \item \textbf{Treating the underlying cause:} optimising blood sugar control in diabetes, or treating hormone deficiency
    \item \textbf{Counselling:} support for couples concerned about fertility or sexual satisfaction
\end{itemize}

\section*{Complications}
\begin{itemize}
    \item Infertility, if the couple wishes to conceive and the semen cannot be recovered
    \item Anxiety and reduced sexual satisfaction
    \item Relationship strain related to fertility concerns
    \item Effects on self-esteem and body image
\end{itemize}

\section*{Prognosis and Outcomes}
Dry orgasm is generally harmless to physical health. The outlook depends on the cause: after prostate surgery, retrograde ejaculation is usually permanent but harmless, while medication-related causes often reverse when the drug is changed. Fertility can frequently be preserved through sperm recovery and assisted reproduction. Most men adjust well once they understand the condition and its implications.

\section*{Prevention and Self-Care}
\begin{itemize}
    \item Discuss ejaculatory changes openly with a doctor, especially after prostate or bladder surgery
    \item Ask about fertility preservation before planned prostate surgery if fathering children is still a goal
    \item Maintain good diabetes control to reduce nerve damage
    \item Seek advice before stopping or changing any medication
    \item Remember that dry orgasm does not usually affect sexual pleasure, which can reduce anxiety
\end{itemize}

\section*{Sources and Further Reading}
The following reputable sources provide further information for healthcare students and patients seeking reliable, current advice about dry orgasm and ejaculatory disorders.
\begin{itemize}
    \item NHS. \textit{Retrograde ejaculation: causes and treatment.} https://www.nhs.uk/conditions/retrograde-ejaculation/
    \item Mayo Clinic. \textit{Retrograde ejaculation: symptoms and causes.} https://www.mayoclinic.org/diseases-conditions/retrograde-ejaculation/
    \item MedlinePlus. \textit{Retrograde ejaculation.} https://medlineplus.gov/retrogradeejaculation.html
    \item MSD Manual. \textit{Ejaculation disorders: professional overview.} https://www.msdmanuals.com/
    \item Urology Care Foundation. \textit{Retrograde ejaculation patient information.} https://www.urologyhealth.org/
    \item Healthy Male Australia. \textit{Male sexual health information.} https://www.healthymale.org.au/
\end{itemize}
